%% =============================================================================
%%  config.tex — caderno de estudo de livro
%% =============================================================================

% ---- identificação do caderno ------------------------------------------------
\newcommand\ebTitulo{Enjoy Life Forever!}
\newcommand\ebSubtitulo{Study notebook, lesson by lesson}
\newcommand\ebIncipit{ἀνακρίνοντες τὰς γραφάς}   % Ac 17:11 “examining the Scriptures”
\newcommand\ebSerie{Bible Study Notebooks}
\newcommand\ebNumero{1}
\newcommand\ebAutor{Isaac Kosloski}
\newcommand\ebData{October 2026}
\newcommand\ebEpigrafe{Those seeking Jehovah will praise him. May you enjoy life forever.}
\newcommand\ebEpigrafeRef{Psalm 22:26 — July 2026}

% ---- o livro estudado ---------------------------------------------------------
\newcommand\livroNome{Enjoy Life Forever!---An Interactive Bible Course}
\newcommand\livroEditora{Watch Tower Bible and Tract Society of Pennsylvania}
\newcommand\livroEdicao{English edition, PDF}
\newcommand\livroDados{dados/livros/lff-en}   % licoes.tsv e textos.tsv
% Nome da unidade: por padrão vem do idioma (Lição / Lesson / Lección).
% Para outro nome, troque aqui: \newcommand\livroUnidade{Capítulo}
\newcommand\livroUnidade{\EBtxt{licao}}
\newcommand\livroUnidadePlural{\EBtxt{licoes}}
\newcommand\livroParte{Section}   % o livro em inglês divide-se em Sections
% impresso: linhas e caixinhas para escrever à mão | digital: só licoes/NN.tex
\newcommand\livroModo{impresso}
\newcommand\livroLinhas{5}
% Lições a incluir: vazio = todas; ex.: 1-12 ou 1,4,7
\newcommand\livroLicoes{}

% ---- aparência ----------------------------------------------------------------
% Idioma da interface (rótulos, tabelas, índices): pt | en | es | …
% Os textos ficam em compartilhado/idiomas/<codigo>.tex
\newcommand\ebIdioma{en}
\newcommand\ebTema{purpura}          % lapis | purpura | oliveira | sepia
\newcommand\ebPapel{a4}         % a4 | b5
\newcommand\ebFundoPagina{nao}
\newcommand\ebAbas{sim}
\newcommand\ebAbasPor{parte}
\newcommand\ebInterlinearModo{estudo}
\newcommand\ebAcentosHebraicos{nao}
\newcommand\ebFontesPersonalizadas{}
